
\documentclass{article}
\usepackage{colortbl}
\usepackage{makecell}
\usepackage{multirow}
\usepackage{supertabular}

\begin{document}

\newcounter{utterance}

\centering \large Interaction Transcript for game `referencegame', experiment `number\_grids', episode 1 with reflexion{-}Qwen3.5{-}9B.
\vspace{24pt}

{ \footnotesize  \setcounter{utterance}{1}
\setlength{\tabcolsep}{0pt}
\begin{supertabular}{c@{$\;$}|p{.15\linewidth}@{}p{.15\linewidth}p{.15\linewidth}p{.15\linewidth}p{.15\linewidth}p{.15\linewidth}}
    \# & $\;$A & \multicolumn{4}{c}{Game Master} & $\;\:$B\\
    \hline

    \theutterance \stepcounter{utterance}  
    & & \multicolumn{4}{p{0.6\linewidth}}{
        \cellcolor[rgb]{0.9,0.9,0.9}{
            \makecell[{{p{\linewidth}}}]{
                \texttt{\tiny{[P1$\langle$GM]}}
                \texttt{You are given three grids, where each of them is 5 by 5 in size.} \\
\texttt{Grids have empty cells marked with "▢" and filled cells marked with "X".} \\
\texttt{Your task is to generate a referring expression that best describes the target grid while distinguishing it from the two other distractor grids.} \\
\texttt{The first grid is the target grid, and the following two grids are the distractors.} \\
\\ 
\texttt{Target grid:} \\
\\ 
\texttt{X X X X □} \\
\texttt{X □ □ □ □} \\
\texttt{X X X X □} \\
\texttt{□ □ □ X □} \\
\texttt{X X X X □} \\
\\ 
\texttt{Distractor grid 1:} \\
\\ 
\texttt{X X X X □} \\
\texttt{X □ □ □ □} \\
\texttt{X X X X □} \\
\texttt{X □ □ □ X} \\
\texttt{X X X X □} \\
\\ 
\texttt{Distractor grid 2:} \\
\\ 
\texttt{□ X X X X} \\
\texttt{□ X □ □ X} \\
\texttt{□ X X X X} \\
\texttt{□ □ □ □ □} \\
\texttt{□ X X X X} \\
\\ 
\texttt{Instruction: Describe the target grid.} \\
\texttt{Generate the referring expression starting with the tag "Expression: " for the given target grid. Omit any other text.} \\
            }
        }
    }
    & \\ \\

    \theutterance \stepcounter{utterance}  
    & \multicolumn{4}{p{0.6\linewidth}}{
        \cellcolor[rgb]{0.9,0.9,0.9}{
            \makecell[{{p{\linewidth}}}]{
                \texttt{\tiny{[P1$\rangle$GM]}}
                \texttt{Expression: The grid with a solid block of X's in the first three rows except for the last column, a row of empty cells in the fourth row except for the fourth cell, and a solid block of X's in the fifth row except for the last column.} \\
            }
        }
    }
    & & \\ \\

    \theutterance \stepcounter{utterance}  
    & & & \multicolumn{2}{p{0.3\linewidth}}{
        \cellcolor[rgb]{0.9,0.9,0.9}{
            \makecell[{{p{\linewidth}}}]{
                \texttt{\tiny{[GM$|$GM]}}
                \texttt{Expression: The grid with a solid block of X's in the first three rows except for the last column, a row of empty cells in the fourth row except for the fourth cell, and a solid block of X's in the fifth row except for the last column.} \\
            }
        }
    }
    & & \\ \\

    \theutterance \stepcounter{utterance}  
    & & \multicolumn{4}{p{0.6\linewidth}}{
        \cellcolor[rgb]{0.9,0.9,0.9}{
            \makecell[{{p{\linewidth}}}]{
                \texttt{\tiny{[GM$\rangle$P2]}}
                \texttt{You are given three grids, where each of them is 5 by 5 in size.} \\
\texttt{Grids have empty cells marked with "▢" and filled cells marked with "X".} \\
\texttt{You are also given a referring expression that describes one of the given grids.} \\
\texttt{Your task is to select the grid that matches the given referring expression.} \\
\texttt{Generate only the number (in text) of the grid that matches the given expression by selecting first, second, or third.} \\
\\ 
\texttt{First:} \\
\\ 
\texttt{X X X X □} \\
\texttt{X □ □ □ □} \\
\texttt{X X X X □} \\
\texttt{X □ □ □ X} \\
\texttt{X X X X □} \\
\\ 
\texttt{Second:} \\
\\ 
\texttt{X X X X □} \\
\texttt{X □ □ □ □} \\
\texttt{X X X X □} \\
\texttt{□ □ □ X □} \\
\texttt{X X X X □} \\
\\ 
\texttt{Third:} \\
\\ 
\texttt{□ X X X X} \\
\texttt{□ X □ □ X} \\
\texttt{□ X X X X} \\
\texttt{□ □ □ □ □} \\
\texttt{□ X X X X} \\
\\ 
\texttt{Expression: The grid with a solid block of X's in the first three rows except for the last column, a row of empty cells in the fourth row except for the fourth cell, and a solid block of X's in the fifth row except for the last column.} \\
\texttt{Question: Which grid does the expression refer to?} \\
\texttt{Start with the tag "Answer: ", followed by your selection. Omit any other text.} \\
            }
        }
    }
    & \\ \\

    \theutterance \stepcounter{utterance}  
    & & & \multicolumn{4}{p{0.6\linewidth}}{
        \cellcolor[rgb]{0.9,0.9,0.9}{
            \makecell[{{p{\linewidth}}}]{
                \texttt{\tiny{[GM$\langle$P2]}}
                \texttt{Answer: First} \\
            }
        }
    }
     \\ \\

    \theutterance \stepcounter{utterance}  
    & & & \multicolumn{2}{p{0.3\linewidth}}{
        \cellcolor[rgb]{0.9,0.9,0.9}{
            \makecell[{{p{\linewidth}}}]{
                \texttt{\tiny{[GM$|$GM]}}
                \texttt{first} \\
            }
        }
    }
    & & \\ \\

    \theutterance \stepcounter{utterance}  
    & \multicolumn{4}{p{0.6\linewidth}}{
        \cellcolor[rgb]{0.95,0.95,0.95}{
            \makecell[{{p{\linewidth}}}]{
                \texttt{\tiny{[P2$\rightarrow$P1]}}
                \texttt{Post{-}episode Reflexion (not sent to GM):} \\
\texttt{New plan: My previous description was too vague and relied on relative terms like "solid block" and "except for," which failed to uniquely distinguish the target from Distractor 1 (which has a single X in the bottom{-}right corner) and Distractor 2 (which has a completely different pattern). Next time, I will construct a description based on precise row{-}by{-}row coordinates and specific cell states. I will explicitly state that Row 4 has exactly one filled cell at column 4, whereas Distractor 1 has two filled cells in Row 4 and Distractor 2 has no filled cells in Row 4. I will also clarify that Rows 1, 3, and 5 are identical across all grids, making the unique feature of Row 4 the sole differentiator, ensuring the expression is both accurate and exclusive to the target.} \\
            }
        }
    }
    & & \\ \\

    \theutterance \stepcounter{utterance}  
    & & & \multicolumn{4}{p{0.6\linewidth}}{
        \cellcolor[rgb]{0.95,0.95,0.95}{
            \makecell[{{p{\linewidth}}}]{
                \texttt{\tiny{[P1$\rightarrow$P2]}}
                \texttt{Post{-}episode Reflexion (not sent to GM):} \\
\texttt{New plan: I incorrectly identified the First grid because I misread the pattern in the fourth row. The expression requires a row of empty cells except for the fourth cell (which is an 'X'), but the First grid has an 'X' in the fourth cell and empty cells elsewhere, matching the description of the Second grid's fourth row pattern (□ □ □ X □). I should have carefully compared the specific cell positions in the fourth row of each grid against the "row of empty cells except for the fourth cell" constraint before making my selection. Next time, I will verify the exact position of the single 'X' in the fourth row for each grid to ensure it aligns with the expression.} \\
            }
        }
    }
     \\ \\

\end{supertabular}
}

\end{document}
